% !TEX root = ../main.tex
% ============================================================
% Chapter 5 — The Frequency Transfer Operator
% Draft focus: background, motivation, operator definitions,
% and the bridge from solution transfer to learned preconditioning.
% ============================================================

\chapter{The frequency transfer operator}
\label{ch:frequency-transfer}

\providecommand{\chfiveincludegraphics}[2]{%
    \IfFileExists{#2}{%
        \includegraphics[width=#1]{#2}%
    }{%
        \fbox{%
            \parbox[c][0.22\textheight][c]{0.86\linewidth}{%
                \centering
                \textbf{Figure placeholder}\\[0.6em]
                \texttt{\detokenize{#2}}%
            }%
        }%
    }%
}

Chapter~\ref{chapter:classical_solvers} showed that the Helmholtz equation
becomes increasingly difficult for classical iterative methods as the
frequency grows. Standard preconditioners improve the linear system at a fixed
frequency, but they do not use the fact that wavefields at neighbouring
frequencies are physically related. This chapter introduces the central idea of
the thesis: use low-frequency information to help solve high-frequency
Helmholtz systems.

The starting point is simple. In a homogeneous medium, solutions at two
frequencies have the same geometric structure but different oscillation rates.
The low-frequency field contains the large-scale propagation pattern. The
high-frequency field contains the same pattern with shorter wavelength and a
different global amplitude scale. This motivates a frequency transfer operator
that maps between solutions at different frequencies.

The chapter has two purposes. First, it explains why such a transfer operator
is plausible from the analytical structure of the Helmholtz Green's function.
Second, it clarifies why a solver cannot blindly use a solution-transfer
network as a preconditioner. A Krylov preconditioner acts on residuals and
returns corrections, not complete solution fields. This distinction is crucial:
solution transfer, residual restriction, and correction prolongation are
related, but they are not the same mathematical object.

% ============================================================
\section{Frequency coherence of Helmholtz fields}
\label{sec:ch5-frequency-coherence}
% ============================================================

Consider the time-harmonic Helmholtz equation in a homogeneous medium with
unit wave speed,
\begin{equation}
    \mathcal{A}_\omega u
    :=
    \left(-\Delta - \omega^2\right)u
    =
    f,
    \label{eq:ch5-helmholtz-continuous}
\end{equation}
where $\omega$ is the angular frequency. In free space, the two-dimensional
outgoing Green's function for a point source at $\mathbf{x}_s$ is
\begin{equation}
    G_\omega(\mathbf{x},\mathbf{x}_s)
    =
    \frac{i}{4} H_0^{(1)}(\omega r),
    \qquad
    r = \|\mathbf{x}-\mathbf{x}_s\|_2,
    \label{eq:ch5-greens-function}
\end{equation}
where $H_0^{(1)}$ is the Hankel function of the first kind. For large argument,
\begin{equation}
    H_0^{(1)}(z)
    \sim
    \sqrt{\frac{2}{\pi z}}
    \exp\!\left(i(z-\pi/4)\right),
    \qquad |z|\rightarrow\infty.
    \label{eq:ch5-hankel-asymptotic}
\end{equation}
Substitution into~\eqref{eq:ch5-greens-function} gives the far-field form
\begin{equation}
    G_\omega(\mathbf{x},\mathbf{x}_s)
    \sim
    C_\omega r^{-1/2} e^{i\omega r},
    \qquad
    C_\omega = \frac{i}{4}\sqrt{\frac{2}{\pi\omega}}e^{-i\pi/4}.
    \label{eq:ch5-green-far-field}
\end{equation}

Equation~\eqref{eq:ch5-green-far-field} separates the wavefield into three
parts. The factor $r^{-1/2}$ is the two-dimensional geometric spreading
envelope. It depends on source--receiver distance but not on frequency. The
constant $C_\omega$ contains the global frequency-dependent amplitude scaling,
with $|C_\omega|\propto \omega^{-1/2}$. The phase factor $e^{i\omega r}$
contains the oscillation rate. Increasing $\omega$ therefore increases the
number of phase cycles along the same propagation path while preserving the
large-scale geometric envelope.

For a frequency pair $(\omega_L,\omega_H)$ with $\omega_H=2\omega_L$, the
far-field amplitude ratio is
\begin{equation}
    \frac{|C_{\omega_H}|}{|C_{\omega_L}|}
    =
    \sqrt{\frac{\omega_L}{\omega_H}}
    =
    \frac{1}{\sqrt{2}}.
    \label{eq:ch5-amplitude-ratio}
\end{equation}
This does not mean that a high-frequency field is a pointwise rescaled
low-frequency field. The phase changes along each propagation path. It does,
however, mean that the two fields are coherent: they share the same underlying
propagation geometry.

\begin{figure}[htbp]
    \centering
    \chfiveincludegraphics{0.92\linewidth}{figures/ch5/green_frequency_coherence.png}
    \caption{Frequency coherence of homogeneous Helmholtz fields. The real
    part of the free-space Green's function is shown for
    $\omega=32,64,128$ with one centred source. The propagation geometry is the
    same in all panels, while the oscillation rate increases with frequency.
    Each panel is normalised by its maximum absolute value to emphasise phase
    structure rather than absolute amplitude.}
    \label{fig:ch5-green-frequency-coherence}
\end{figure}

\begin{figure}[htbp]
    \centering
    \chfiveincludegraphics{0.78\linewidth}{figures/ch5/radial_green_cross_section.png}
    \caption{Radial cross-section of the free-space Green's function for the
    pair $(\omega_L,\omega_H)=(64,128)$. Top: the high-frequency amplitude
    envelope is lower by the expected far-field factor
    $\sqrt{\omega_L/\omega_H}=1/\sqrt{2}$. Bottom: the high-frequency real
    part completes approximately twice as many oscillations over the same
    distance. The shaded near-source region indicates where the far-field
    asymptotic approximation is least accurate.}
    \label{fig:ch5-radial-green-cross-section}
\end{figure}

% ============================================================
\section{Solution transfer}
\label{sec:ch5-solution-transfer}
% ============================================================

Let $u_L$ and $u_H$ solve
\begin{equation}
    \mathcal{A}_L u_L = f,
    \qquad
    \mathcal{A}_H u_H = f,
    \label{eq:ch5-two-frequency-solves}
\end{equation}
where $\mathcal{A}_L=\mathcal{A}_{\omega_L}$ and
$\mathcal{A}_H=\mathcal{A}_{\omega_H}$. The exact solution-transfer operator
from low to high frequency is
\begin{equation}
    \mathcal{T}_{\uparrow}
    :=
    \mathcal{A}_H^{-1}\mathcal{A}_L.
    \label{eq:ch5-exact-solution-transfer-up}
\end{equation}
Indeed, if $u_L=\mathcal{A}_L^{-1}f$, then
\begin{equation}
    \mathcal{T}_{\uparrow}u_L
    =
    \mathcal{A}_H^{-1}\mathcal{A}_L u_L
    =
    \mathcal{A}_H^{-1}f
    =
    u_H.
    \label{eq:ch5-transfer-maps-solutions}
\end{equation}
Similarly, the exact high-to-low solution-transfer operator is
\begin{equation}
    \mathcal{T}_{\downarrow}^{\mathrm{sol}}
    :=
    \mathcal{A}_L^{-1}\mathcal{A}_H.
    \label{eq:ch5-exact-solution-transfer-down}
\end{equation}

This notation is deliberately precise. The identity
\begin{equation}
    \mathcal{A}_H\mathcal{T}_{\uparrow}
    =
    \mathcal{A}_L
    \label{eq:ch5-solution-intertwining}
\end{equation}
holds for the operator~\eqref{eq:ch5-exact-solution-transfer-up}. It says that
applying the high-frequency operator after exact transfer recovers the
low-frequency forcing represented by the original low-frequency solution. This
is the relevant algebraic identity for solution transfer.

One should not confuse~\eqref{eq:ch5-solution-intertwining} with a generic
commutation relation of the form
$\mathcal{A}_H\mathcal{T}=\mathcal{T}\mathcal{A}_L$. The latter would require
the same operator to transfer both fields and right-hand sides in a compatible
way. For solver design, this distinction matters because Krylov methods do not
only manipulate solution fields. They repeatedly manipulate residuals.

In a constant-coefficient periodic or free-space setting, the Fourier transform
diagonalises the Helmholtz operator. Formally,
\begin{equation}
    \widehat{\mathcal{A}_\omega u}(\mathbf{k})
    =
    \left(|\mathbf{k}|^2-\omega^2\right)\hat{u}(\mathbf{k}).
    \label{eq:ch5-fourier-symbol-A}
\end{equation}
The solution-transfer operator~\eqref{eq:ch5-exact-solution-transfer-up}
therefore has Fourier multiplier
\begin{equation}
    \widehat{\mathcal{T}_{\uparrow}}(\mathbf{k})
    =
    \frac{|\mathbf{k}|^2-\omega_L^2}
         {|\mathbf{k}|^2-\omega_H^2},
    \label{eq:ch5-fourier-symbol-transfer}
\end{equation}
with the usual outgoing or PML regularisation near singular frequencies. In
physical space this corresponds to a convolution operator. This explains why a
convolutional neural network is a natural approximation class in homogeneous
media: the ideal map is translation-invariant away from boundaries.

\begin{figure}[htbp]
    \centering
    \chfiveincludegraphics{0.86\linewidth}{figures/ch5/solution_transfer_operator_schematic.png}
    \caption{Exact solution transfer between two frequencies. The same
    right-hand side $f$ defines both $u_L=\mathcal{A}_L^{-1}f$ and
    $u_H=\mathcal{A}_H^{-1}f$. The solution-transfer operator
    $\mathcal{T}_{\uparrow}=\mathcal{A}_H^{-1}\mathcal{A}_L$ maps $u_L$ to
    $u_H$, and satisfies
    $\mathcal{A}_H\mathcal{T}_{\uparrow}u_L=\mathcal{A}_L u_L=f$ on this
    solution manifold.}
    \label{fig:ch5-solution-transfer-schematic}
\end{figure}

% ============================================================
\section{From solution transfer to preconditioning}
\label{sec:ch5-transfer-to-preconditioning}
% ============================================================

A trained network that approximates
\begin{equation}
    T_{\uparrow,\theta}: u_L \mapsto u_H
    \label{eq:ch5-learned-solution-transfer}
\end{equation}
can be useful as a warm start: solve the cheaper low-frequency problem, transfer
the result, and use the prediction as an initial guess for the high-frequency
Krylov solve. If
\begin{equation}
    x_0 = T_{\uparrow,\theta}u_L,
    \label{eq:ch5-warm-start}
\end{equation}
then the initial high-frequency residual is
\begin{equation}
    r_0
    =
    f-\mathbf{A}_H x_0
    =
    \mathbf{A}_H(u_H-x_0).
    \label{eq:ch5-residual-error-relation}
\end{equation}
This equation already reveals a potential failure mode. A small field error
$u_H-x_0$ does not necessarily imply a small residual. Components of the error
in directions where $\mathbf{A}_H$ has large amplification can dominate the
residual even if they are visually small in the field.

This is the conceptual bridge to the later one-dimensional spectral analysis.
The goal is not only to predict a high-frequency field that looks accurate. The
goal is to produce an initial residual, or a preconditioned residual, that is
favourable for FGMRES.

A stronger use of frequency transfer is therefore to apply it inside the
preconditioner at every Krylov iteration. At iteration $j$, FGMRES supplies a
residual $r_H^{(j)}$. A preconditioner must return an approximate correction
$z_H^{(j)}$ satisfying
\begin{equation}
    \mathbf{A}_H z_H^{(j)}
    \approx
    r_H^{(j)}.
    \label{eq:ch5-preconditioner-task}
\end{equation}
This is a different task from predicting $u_H$ from $u_L$.

The learned V-cycle considered in this thesis has the form
\begin{equation}
    \mathbf{M}_{\theta}^{-1}
    =
    \mathcal{P}_{\theta}
    \mathbf{A}_L^{-1}
    \mathcal{R}_{\theta},
    \label{eq:ch5-learned-vcycle}
\end{equation}
where $\mathcal{R}_{\theta}$ maps a high-frequency residual to a low-frequency
right-hand-side or correction representation, $\mathbf{A}_L^{-1}$ applies the
low-frequency solve, and $\mathcal{P}_{\theta}$ maps the resulting correction
back to the high-frequency grid. This structure is analogous to a multigrid
V-cycle:
\begin{equation}
    \text{restriction}
    \;\longrightarrow\;
    \text{coarse solve}
    \;\longrightarrow\;
    \text{prolongation}.
    \label{eq:ch5-vcycle-words}
\end{equation}

\begin{figure}[htbp]
    \centering
    \chfiveincludegraphics{0.80\linewidth}{figures/ch5/learned_vcycle_schematic.png}
    \caption{Learned V-cycle preconditioner applied inside one FGMRES
    iteration. The high-frequency residual $r_H^{(j)}$ is restricted by a
    learned map $\mathcal{R}_{\theta}$, solved on the low-frequency operator
    $\mathbf{A}_L^{-1}$, and prolonged by a learned map
    $\mathcal{P}_{\theta}$ to produce the high-frequency correction
    $z_H^{(j)}=\mathbf{M}_{\theta}^{-1}r_H^{(j)}$.}
    \label{fig:ch5-learned-vcycle-schematic}
\end{figure}

The analogy is useful, but it should not hide the main difficulty. Classical
restriction and prolongation are defined by grid geometry. Here the maps must
also respect frequency-dependent wave physics. In addition, the input to the
preconditioner is a residual, while many neural transfer models are naturally
trained on solution pairs. The thesis therefore distinguishes three objects:
\begin{align}
    T_{\uparrow}^{\mathrm{sol}} &: u_L \mapsto u_H,
    &&\text{solution transfer}, \label{eq:ch5-three-operators-a}\\
    R &: r_H \mapsto r_L \text{ or } e_L,
    &&\text{residual restriction}, \label{eq:ch5-three-operators-b}\\
    P &: e_L \mapsto e_H,
    &&\text{correction prolongation}. \label{eq:ch5-three-operators-c}
\end{align}
Treating these as interchangeable is tempting, but it is not mathematically
justified. Later results show that this distinction is one of the central
practical issues in learned Helmholtz preconditioning.

\begin{figure}[htbp]
    \centering
    \chfiveincludegraphics{0.92\linewidth}{figures/ch5/three_transfer_objects.png}
    \caption{Three transfer objects that must be distinguished. Solution
    transfer maps one complete wavefield to another. Residual restriction maps
    the current high-frequency residual to a low-frequency representation.
    Correction prolongation maps a low-frequency correction back to the
    high-frequency system. These objects have different mathematical roles and
    should not be treated as interchangeable.}
    \label{fig:ch5-three-transfer-objects}
\end{figure}

% ============================================================
\section{Residual amplification and modal safety}
\label{sec:ch5-residual-amplification}
% ============================================================

The simplest setting in which the difference between field accuracy and solver
usefulness becomes visible is the one-dimensional Dirichlet problem. Let
$\mathbf{A}_H$ have eigenpairs $(\lambda_k,v_k)$ with
$\|v_k\|_2=1$. If the initial error is expanded as
\begin{equation}
    e_0 = u_H-x_0 = \sum_k c_k(e_0)v_k,
    \label{eq:ch5-error-expansion}
\end{equation}
then the residual is
\begin{equation}
    r_0 = \mathbf{A}_H e_0
    =
    \sum_k \lambda_k c_k(e_0)v_k.
    \label{eq:ch5-residual-expansion}
\end{equation}
Thus each residual coefficient is the corresponding error coefficient
multiplied by $\lambda_k$:
\begin{equation}
    c_k(r_0)=\lambda_k c_k(e_0).
    \label{eq:ch5-modal-amplification}
\end{equation}

Equation~\eqref{eq:ch5-modal-amplification} is the core warning. A learned
field can have small relative $\ell^2$ error while still producing a large
initial residual if it places small errors in high-$|\lambda_k|$ modes. This is
not a neural-network artefact; it is an algebraic property of the Helmholtz
operator.

For this reason, a learned correction should be judged not only by field error
but also by residual error, preconditioned residual error, and modal
distribution. A safe learned correction is one that improves the residual in
the modes where it is used. This motivates spectral filtering or gating:
compare a learned correction against a trusted baseline, such as a CSL
preconditioner, and accept the learned component only where it reduces the
residual.

The one-dimensional Dirichlet case is especially valuable because the
eigenpairs are analytical. The eigenvectors are sine modes,
\begin{equation}
    v_k(j)
    =
    \sqrt{\frac{2}{N+1}}
    \sin\!\left(\frac{jk\pi}{N+1}\right),
    \qquad
    j,k=1,\dots,N,
    \label{eq:ch5-dirichlet-eigenvectors}
\end{equation}
and the eigenvalues are
\begin{equation}
    \lambda_k(\omega)
    =
    \frac{4}{h^2}
    \sin^2\!\left(\frac{k\pi}{2(N+1)}\right)
    -
    \omega^2.
    \label{eq:ch5-dirichlet-eigenvalues}
\end{equation}
This provides a controlled microscope for checking whether learned transfer
operators help or hurt the solver.

\begin{figure}[htbp]
    \centering
    \chfiveincludegraphics{0.86\linewidth}{figures/ch5/modal_amplification_cartoon.png}
    \caption{Residual amplification in the Dirichlet eigenbasis. A field error
    with small high-mode coefficients can still create a large residual because
    each residual coefficient is multiplied by the corresponding eigenvalue:
    $c_k(r_0)=\lambda_k c_k(e_0)$. This is why field accuracy alone is not a
    sufficient criterion for Krylov acceleration.}
    \label{fig:ch5-modal-amplification-cartoon}
\end{figure}

% ============================================================
\section{Learning objectives}
\label{sec:ch5-learning-objectives}
% ============================================================

The most direct training objective for solution transfer is the complex
relative $\ell^2$ loss
\begin{equation}
    \mathcal{L}_{\mathrm{field}}
    =
    \frac{\|\hat{u}_H-u_H\|_2^2}
         {\|u_H\|_2^2+\varepsilon}.
    \label{eq:ch5-field-loss}
\end{equation}
This objective is natural for warm-start prediction. It is scale-invariant and
penalises errors in the complex field directly. In the experiments, real and
imaginary parts are represented as separate channels but interpreted as one
complex vector.

For preconditioning, a more solver-facing objective is residual-based. If a
network predicts a correction $\hat{e}$ for a residual $r$, then the relevant
loss is
\begin{equation}
    \mathcal{L}_{\mathrm{res}}
    =
    \frac{\|\mathbf{A}\hat{e}-r\|_2^2}
         {\|r\|_2^2+\varepsilon}.
    \label{eq:ch5-residual-loss}
\end{equation}
This loss directly asks whether the correction solves the linearised task
needed by the preconditioner. It is more aligned with FGMRES than
\eqref{eq:ch5-field-loss}, but it is also more difficult: applying
$\mathbf{A}$ amplifies high-frequency components and makes the optimisation
more sensitive to modal errors.

A third possibility is to train or evaluate in a preconditioned residual norm,
for example
\begin{equation}
    \mathcal{L}_{\mathrm{prec}}
    =
    \frac{\|\mathbf{P}^{-1}(\mathbf{A}\hat{e}-r)\|_2^2}
         {\|\mathbf{P}^{-1}r\|_2^2+\varepsilon},
    \label{eq:ch5-preconditioned-residual-loss}
\end{equation}
where $\mathbf{P}^{-1}$ is a trusted baseline preconditioner such as complex
shifted Laplace. This objective measures errors in the norm most closely
related to the actual preconditioned Krylov process.

These objectives answer different questions:
\begin{table}[htbp]
    \centering
    \caption{Different learning objectives and the question each one answers.}
    \label{tab:ch5-objectives}
    \begin{tabular}{@{}p{0.25\linewidth}p{0.35\linewidth}p{0.30\linewidth}@{}}
        \toprule
        Objective & Question answered & Best use \\
        \midrule
        Field loss
        & Does the predicted wavefield match the target solution?
        & Warm starts and solution transfer \\
        Residual loss
        & Does the predicted correction solve $\mathbf{A}e=r$?
        & Correction operators and preconditioning \\
        Preconditioned residual loss
        & Does the correction help after the baseline preconditioner?
        & Solver-facing learned corrections \\
        Modal/gated diagnostics
        & In which spectral components is the learned correction safe?
        & Interpretation and robust hybrid methods \\
        \bottomrule
    \end{tabular}
\end{table}

The later experimental chapters use this distinction to organise the results.
The field-trained models are evaluated not only by relative field error but
also by their residuals and FGMRES iteration counts. Learned V-cycle variants
are evaluated by whether they improve the preconditioned residual at every
iteration, not merely by whether their output resembles a solution field.

% ============================================================
\section{Architectural implications}
\label{sec:ch5-architecture-implications}
% ============================================================

The analytical discussion suggests several architectural requirements.

First, the model needs a non-trivial receptive field. The phase transformation
from $\omega_L$ to $\omega_H$ depends on propagation distance and cannot be
represented by independent pointwise scaling. Convolutional architectures are
therefore natural because they can learn local integral kernels.

Second, the representation must preserve complex phase. The real and imaginary
parts are kept as separate channels. An amplitude--phase representation would
introduce artificial discontinuities at phase wraps.

Third, the model should have a mechanism for multi-scale information. Low
frequencies carry broad wave envelopes, while high frequencies require
oscillatory detail. A U-Net is a natural choice because the encoder aggregates
large-scale context and the decoder reconstructs fine-scale structure through
skip connections.

Fourth, the architecture alone does not guarantee solver usefulness. Even a
network with excellent field accuracy may create residual-unfriendly modes.
Therefore the solver-facing evaluation must include residual norms, modal
diagnostics, and FGMRES convergence.

The thesis focuses on U-Net transfer models for this reason. The U-Net is
expressive enough to learn solution transfer, while the later chapters test
whether its outputs can be made safe and useful inside a Krylov preconditioner.

\begin{figure}[htbp]
    \centering
    \chfiveincludegraphics{0.90\linewidth}{figures/ch5/unet_frequency_transfer_schematic.png}
    \caption{U-Net as a frequency-transfer model. The encoder aggregates
    large-scale wavefield structure, while skip connections preserve
    oscillatory detail needed by the decoder. The same architectural family can
    be used for solution transfer or correction prediction, but the target and
    loss determine which mathematical object is learned.}
    \label{fig:ch5-unet-frequency-transfer-schematic}
\end{figure}

% ============================================================
\section{Summary}
\label{sec:ch5-summary}
% ============================================================

This chapter introduced frequency transfer as a way to exploit the physical
coherence between Helmholtz solutions at neighbouring frequencies. In the
homogeneous setting, the Green's function separates into a shared geometric
envelope, a frequency-dependent amplitude scale, and a frequency-dependent
phase. This makes solution transfer from $\omega_L$ to $\omega_H$ a plausible
learning problem.

The chapter also clarified the central solver issue. A neural operator trained
to map $u_L$ to $u_H$ is a solution-transfer operator. A preconditioner,
however, acts on residuals and returns corrections. Because residuals amplify
field errors through the Helmholtz operator, field accuracy alone is not a
sufficient criterion for Krylov acceleration. This motivates the later
experiments: first diagnose learned transfer in the analytically controlled 1D
Dirichlet setting, then test whether spectrally safe learned corrections can be
used inside FGMRES.

% ============================================================
% Suggested visual plan for this chapter
% ============================================================
%
% 1. figures/ch5/green_frequency_coherence.png
%    Real parts of Green's functions at omega=32,64,128.
%    Purpose: show same geometry, increasing oscillation rate.
%
% 2. figures/ch5/radial_green_cross_section.png
%    Radial amplitude and real-part cross sections for one frequency pair.
%    Purpose: show 1/sqrt(2) amplitude scaling and doubled phase cycles.
%
% 3. figures/ch5/solution_transfer_operator_schematic.png
%    Diagram of f -> A_L^{-1} -> u_L -> T_up -> u_H and f -> A_H^{-1} -> u_H.
%    Purpose: define solution transfer cleanly.
%
% 4. figures/ch5/three_transfer_objects.png
%    Three separate rows: solution transfer, residual restriction,
%    correction prolongation.
%    Purpose: prevent the most important conceptual confusion.
%
% 5. figures/ch5/modal_amplification_cartoon.png
%    Small high-mode field error multiplied by large lambda becomes large
%    residual.
%    Purpose: foreshadow the 1D spectral results chapter.
%
% 6. figures/ch5/learned_vcycle_schematic.png
%    r_H -> R_theta -> A_L^{-1} -> P_theta -> z_H inside FGMRES.
%    Purpose: show the actual preconditioner used later.
%
% 7. figures/ch5/unet_frequency_transfer_schematic.png
%    U-Net as a transfer architecture, with encoder/decoder and skip paths.
%    Purpose: support the architecture implications section.
